\section{Picard--Fuchs forcing}\label{sec:forcing}

The locations and orders of the stationary circles are controlled by an
inhomogeneous period equation.  Its forcing depends on the marked section
$P$, so the moving endpoint must be retained throughout the reduction.
In the curve \eqref{eq:weierstrass}, set
\[
\begin{gathered}
s=h^2-4T,\qquad X=u+s/12,\qquad Y=v+(hu-T)/2,\\
Y^2=f(X,h)=X^3+aX+b,\\
a=-\frac{s^2+24hT}{48},\qquad
b=\frac{s^3+36hTs+216T^2}{864}.
\end{gathered}
\]
The marked point is $(X_P,Y_P)=(s/12,T/2)$, and the relative differential
is still $\omega=\dd X/(2Y)$: no energy-dependent rescaling has been made.
Write $q_0=8h-9$ and $H=32T+3h$, and retain the discriminant
polynomial $\delta$ from \eqref{eq:critical}.

\begin{proposition}\label{prop:pf}
On any sufficiently small complex parameter domain with $q_0\delta\ne0$,
every period $L$ of $\omega$ and every local integral
$I=\int_O^P\omega$ satisfy
\begin{align}
\mathcal L
&=\partial_h^2+\left(\frac{\delta'}{\delta}-\frac8{q_0}\right)
  \partial_h+\frac{8h^3-18h^2+9h-12T}{q_0\delta},
\label{eq:pf-operator}\\
\mathcal L L&=0,\qquad \mathcal L I=-\frac{H}{q_0\delta}.
\label{eq:pf-forcing}
\end{align}
For a real positive period $L$ and the real positive arc integral $I_j$
from $O$ to $jP$, put $\rho=I_j/L$.  Here $j=1$ on the lower and middle
intervals, while $j=2$ on the upper interval.  Then
\begin{equation}\label{eq:wronskian}
J_\rho=\frac{\delta}{q_0}L^2\rho',\qquad
J_\rho'=-\frac{jHL}{q_0^2}.
\end{equation}
These identities use the actual circle return: $+P$ below the saddle and
$+2P$ above it.
\end{proposition}

\begin{proof}
\emph{Marked reduction.}
In this proof $D=\partial_h$ acts at fixed $X$ and $T$, whereas primes on
the marked coordinates denote their total derivatives.  Thus
$DY=f_h/(2Y)$, and $f_h(P)$ means differentiation before evaluation at
$X_P$.  With $\omega_1=X\omega$, define
\[
\begin{gathered}
\alpha_{\mathrm{PF}}=-\frac{\delta'}{12\delta},\qquad
\beta_{\mathrm{PF}}=-\frac{q_0}{\delta},\qquad
\gamma_{\mathrm{PF}}=\frac{a\beta_{\mathrm{PF}}}{3},\\
\kappa_{\mathrm{PF}}=\frac{\beta_{\mathrm{PF}}'}{\beta_{\mathrm{PF}}}
=\frac8{q_0}-\frac{\delta'}{\delta},\\
V_{\mathrm{PF}}=\alpha_{\mathrm{PF}}'+\alpha_{\mathrm{PF}}^2
+\beta_{\mathrm{PF}}\gamma_{\mathrm{PF}}
-\alpha_{\mathrm{PF}}\kappa_{\mathrm{PF}},\\
r_0=2\alpha_{\mathrm{PF}}X-2\beta_{\mathrm{PF}}X^2
-\tfrac43a\beta_{\mathrm{PF}},\\
r_1=2\alpha_{\mathrm{PF}}X^2+\tfrac23a\beta_{\mathrm{PF}}X
+2b\beta_{\mathrm{PF}},\qquad R_i=\frac{r_i}{2Y}.
\end{gathered}
\]
Differentiating the displayed coefficients $a,b$ gives
$a'=-4a\alpha_{\mathrm{PF}}+6b\beta_{\mathrm{PF}}$ and
$b'=-6b\alpha_{\mathrm{PF}}-4a^2\beta_{\mathrm{PF}}/3$.
Equating coefficients of $X$ then gives the two polynomial identities
\begin{equation}\label{pf:reduction-polynomials}
\begin{aligned}
-f_h&=2(\alpha_{\mathrm{PF}}+\beta_{\mathrm{PF}}X)f
       +2f(r_0)_X-f_Xr_0,\\
-Xf_h&=2(\gamma_{\mathrm{PF}}-\alpha_{\mathrm{PF}}X)f
       +2f(r_1)_X-f_Xr_1.
\end{aligned}
\end{equation}
Multiplication by $\dd X/(4Y^3)$ proves the exact differential identities
\begin{equation}\label{pf:gauss-manin}
\begin{aligned}
D\omega&=\alpha_{\mathrm{PF}}\omega
+\beta_{\mathrm{PF}}\omega_1+\dd_XR_0,\\
D\omega_1&=\gamma_{\mathrm{PF}}\omega
-\alpha_{\mathrm{PF}}\omega_1+\dd_XR_1.
\end{aligned}
\end{equation}
Here $\dd_X$ denotes the relative differential.  The same polynomials
also give $R_1-XR_0=\beta_{\mathrm{PF}}Y$.
Differentiating the first line of \eqref{pf:gauss-manin} and eliminating
$\omega_1$ yields
\begin{equation}\label{pf:exact-primitive}
\begin{aligned}
(D^2-\kappa_{\mathrm{PF}}D-V_{\mathrm{PF}})\omega
&=\dd_XQ_{\mathrm{PF}},\\
Q_{\mathrm{PF}}
&=(D+\alpha_{\mathrm{PF}}-\kappa_{\mathrm{PF}})R_0
  +\beta_{\mathrm{PF}}R_1\\
&=\frac{2V_{\mathrm{PF}}X-\beta_{\mathrm{PF}}a'}{2Y}
  -\frac{r_0f_h}{4Y^3}.
\end{aligned}
\end{equation}
For the last equality, the numerator identity is
$Dr_0+(\alpha_{\mathrm{PF}}-\kappa_{\mathrm{PF}})r_0
+\beta_{\mathrm{PF}}r_1=2V_{\mathrm{PF}}X-\beta_{\mathrm{PF}}a'$.
Substitution and a common denominator give
$-V_{\mathrm{PF}}=(8h^3-18h^2+9h-12T)/(q_0\delta)$;
hence the operator in \eqref{pf:exact-primitive} is \eqref{eq:pf-operator}.
Integrating the exact differential over any locally transported closed
cycle proves $\mathcal L L=0$.

\emph{The two endpoints.}
Near $O$ choose the local coordinate $t=X^{-1/2}$ with
$Y=t^{-3}(1+at^4+bt^6)^{1/2}$.  Then
$\omega=-(1+O(t^4))\dd t$, so $I$ converges, but the individual primitives
have poles:
\[
R_0=-\frac{\beta_{\mathrm{PF}}}{t}
  +\alpha_{\mathrm{PF}}t+O(t^3),\qquad
R_1=\frac{\alpha_{\mathrm{PF}}}{t}
  +\gamma_{\mathrm{PF}}t+O(t^3).
\]
Their combination in \eqref{pf:exact-primitive} has pole coefficient
$-\beta_{\mathrm{PF}}'
-(\alpha_{\mathrm{PF}}-\kappa_{\mathrm{PF}})\beta_{\mathrm{PF}}
+\beta_{\mathrm{PF}}\alpha_{\mathrm{PF}}=0$.
There is no constant term, so $Q_{\mathrm{PF}}=O(t)$ and
$Q_{\mathrm{PF}}(O)=0$.  In particular, the two divergent primitives
have not been separately assigned values at $O$.

Apply Leibniz' rule twice, first with a fixed $X$ cutoff near $O$ and
then let the cutoff tend to $O$.  The preceding expansions control this
limit and give all moving-endpoint terms:
\begin{equation}\label{pf:leibniz}
\begin{aligned}
\mathcal L I&=Q_{\mathrm{PF}}(P)-Q_{\mathrm{PF}}(O)
+\frac{X_P''-\kappa_{\mathrm{PF}}X_P'}{2Y_P}
\\
&\hspace{1em}-\frac{2X_P'f_h(P)+(X_P')^2f_X(P)}{4Y_P^3}.
\end{aligned}
\end{equation}
Although $Y_P'=0$, the identity is
$f_h(P)+f_X(P)X_P'=0$, not $f_h(P)=0$.
In fact $X_P'=h/6$, $X_P''=1/6$, $f_X(P)=-Th/2$ and
$f_h(P)=Th^2/12$.  Thus the moving terms in \eqref{pf:leibniz} sum to
$-h^3/(36T^2)+(1-\kappa_{\mathrm{PF}}h)/(6T)$, while evaluation of
\eqref{pf:exact-primitive} gives
\[
Q_{\mathrm{PF}}(P)=\frac{h^3}{36T^2}
+\frac{\kappa_{\mathrm{PF}}h-1}{6T}-\frac{H}{q_0\delta}.
\]
The cancellation proves the inhomogeneous identity
\eqref{eq:pf-forcing}.

\emph{The real return and its Wronskian.}
For $j=1$ the positive real arc is a choice of $I$.  For $j=2$, a local
complex elliptic logarithm satisfies
$\log(2P)=2\log(P)$ modulo the period lattice.  Consequently the actual
real arc $I_2$ differs from $2I$ by a locally fixed integral combination
of periods, all annihilated by $\mathcal L$.  This proves
$\mathcal L I_j=-jH/(q_0\delta)$ in both cases; it does not identify a
complex arc to $P$ on the upper interval with a real single-circle step.
For $W=L I_j'-I_jL'=L^2\rho'$, subtraction of the two scalar equations
gives
\[
W'=\kappa_{\mathrm{PF}}W-\frac{jHL}{q_0\delta}.
\]
Since $(\delta/q_0)'/(\delta/q_0)=-\kappa_{\mathrm{PF}}$,
multiplication by $\delta/q_0$ proves \eqref{eq:wronskian}.
\end{proof}

The zero of $q_0$ is introduced by eliminating $\omega_1$; it is not a
singular fibre.  The real rotation function remains analytic there.
The next section treats this apparent pole and fixes the integration
constants in \eqref{eq:wronskian} using the actual finite and infinite
ends of the real circles.
